% Part: first-order-logic
% Chapter: sequent-calculus
% Section: derivations

\documentclass[../../../include/open-logic-section]{subfiles}

\begin{document}

\iftag{FOL}
      {\olfileid{fol}{seq}{der}}
      {\olfileid{pl}{seq}{der}}

\olsection{\usetoken{P}{derivation}}

\begin{explain}
سبق أن بيّنا صورة التتابعية الابتدائية، وقدمنا قواعد الاستدلال.
وتتولد ال!!^{derivation}s في حساب التتابعيات استقرائيًا من هذه
المعطيات: فكل !!{derivation} إما أن يكون تتابعية ابتدائية قائمة
بنفسها، وإما أن يتألف من واحد أو اثنين من ال!!{derivation}s يتلوهما
استدلال.
\end{explain}

\begin{defn}[$\Log{LK}$: !!{derivation}]
إن \emph{$\Log{LK}$-!!{derivation}} لتتابعية~$S$ شجرة منتهية من
التتابعيات تحقق الشروط الآتية:
\begin{enumerate}
\item التتابعيات العليا في الشجرة تتابعيات ابتدائية.
\item التتابعية السفلى في الشجرة هي~$S$.
\item كل تتابعية في الشجرة عدا~$S$ مقدمة لتطبيق صحيح لقاعدة استدلال
  تقع نتيجته تحت تلك التتابعية مباشرةً في الشجرة.
\end{enumerate}
ونسمّي عندئذٍ~$S$ \emph{التتابعية النهائية} لل!!{derivation}، ونقول
إن~$S$ \emph{!!{derivable} في $\Log{LK}$} (أو
$\Log{LK}$-!!{derivable}).
\end{defn}

\begin{ex}
كل تتابعية ابتدائية، مثل~$!C \Sequent !C$، هي !!a{derivation}. ويمكننا
أن نحصل منها على !!{derivation} جديد بتطبيق القاعدة
$\LeftR{\Weakening}$ مثلًا،
\begin{prooftree}
\Axiom$ \Gamma \fCenter \Delta $
\RightLabel{\LeftR{\Weakening}}
\UnaryInf$ !A, \Gamma \fCenter \Delta$
\end{prooftree}
غير أن المقصود بالقاعدة أن تكون عامة: فيمكننا أن نستبدل~$!A$ في
القاعدة بأي !!{sentence}، ومن ذلك~$!D$ مثلًا. ولكي تطابق المقدمة
تتابعيتنا الابتدائية~$!C \Sequent !C$، لا بد أن يكون كل من~$\Gamma$
و$\Delta$ هو~$!C$ فحسب، وعندئذ تكون النتيجة
$!D, !C \Sequent !C$. ومن ثم فإن ما يأتي !!a{derivation}:
\begin{prooftree}
\Axiom$ !C \fCenter !C $
\RightLabel{\LeftR{\Weakening}}
\UnaryInf$ !D, !C \fCenter !C$
\end{prooftree}
يمكننا الآن تطبيق قاعدة أخرى، ولتكن~$\LeftR{\Exchange}$، وهي تتيح لنا
تبديل !!{sentence}s اثنتين في الجانب الأيسر. ومن ثم فإن ما يأتي أيضًا
!!{derivation} صحيح:
\begin{prooftree}
\Axiom$ !C \fCenter !C $
\RightLabel{\LeftR{\Weakening}}
\UnaryInf$ !D, !C \fCenter !C$
\RightLabel{\LeftR{\Exchange}}
\UnaryInf$ !C, !D \fCenter !C$
\end{prooftree}
وفي هذا التطبيق للقاعدة، التي عُرضت على الصورة
\begin{prooftree}
\Axiom$ \Gamma, !A, !B, \Pi \fCenter \Delta $
\RightLabel{\LeftR{\Exchange}}
\UnaryInf$ \Gamma, !B, !A, \Pi \fCenter \Delta,$
\end{prooftree}
كان كل من~$\Gamma$ و$\Pi$ خاليًا، وكانت~$\Delta$ هي~$!C$، وأدى
دوري~$!A$ و$!B$ كل من~$!D$ و~$!C$ على الترتيب. وبالطريقة نفسها
تقريبًا نرى أيضًا أن
\begin{prooftree}
\Axiom$ !D \fCenter !D $
\RightLabel{\LeftR{\Weakening}}
\UnaryInf$ !C, !D \fCenter !D$
\end{prooftree}
هو !!a{derivation}. والآن يمكننا أخذ هذين ال!!{derivation}s ودمجهما
باستعمال~$\RightR{\land}$. وقد عُرضت تلك القاعدة على الصورة
\begin{prooftree}
\Axiom$\Gamma \fCenter \Delta, !A$
\Axiom$ \Gamma \fCenter \Delta, !B$
\RightLabel{\RightR{\land}}
\BinaryInf$ \Gamma \fCenter \Delta, !A \land !B $
\end{prooftree}
وفي حالتنا، يجب أن تطابق المقدمتان التتابعيتين الأخيرتين
لل!!{derivation}s المنتهيين بالمقدمتين. وهذا يعني أن~$\Gamma$ هي
$!C, !D$، وأن~$\Delta$ خالية، وأن~$!A$ هي~$!C$ و$!B$ هي~$!D$.
ومن ثم تكون النتيجة، لكي يكون الاستدلال صحيحًا،
$!C, !D \Sequent !C \land !D$.
\begin{prooftree}
\Axiom$ !C \fCenter !C $
\RightLabel{\LeftR{\Weakening}}
\UnaryInf$ !D, !C \fCenter !C$
\RightLabel{\LeftR{\Exchange}}
\UnaryInf$ !C, !D \fCenter !C$
\Axiom$ !D \fCenter !D $
\RightLabel{\LeftR{\Weakening}}
\UnaryInf$ !C, !D \fCenter !D$
\RightLabel{\RightR{\land}}
\BinaryInf$ !C, !D \fCenter !C \land !D $
\end{prooftree}
وبالطبع يمكننا أيضًا عكس المقدمتين؛ وعندئذ تكون~$!A$ هي~$!D$،
وتكون~$!B$ هي~$!C$. 
\begin{prooftree}
\Axiom$ !D \fCenter !D $
\RightLabel{\LeftR{\Weakening}}
\UnaryInf$ !C, !D \fCenter !D$
\Axiom$ !C \fCenter !C $
\RightLabel{\LeftR{\Weakening}}
\UnaryInf$ !D, !C \fCenter !C$
\RightLabel{\LeftR{\Exchange}}
\UnaryInf$ !C, !D \fCenter !C$
\RightLabel{\RightR{\land}}
\BinaryInf$ !C, !D \fCenter !D \land !C $
\end{prooftree}
\end{ex}

\end{document}
